\section{Random series of independent samples}

We now look at the related topic of infinite series
of independent r.v.'s. When can such a series be said
to converge? A key technical result that will help us
answer this question in some cases is the following
beautiful inequality due to Kolmogorov.

\begin{thm}[Kolmogorov's maximal inequality]
Assume that $X_1,X_2,\ldots,X_n$ are independent
r.v.'s with finite variances, and let
$S_k=\sum_{j=1}^k X_j$. Then

\[
\mathrm{P}\left( \max_{1\le k\le n} |S_k-\mathbf{E} (S_k)| > t \right) \le \frac{\mathbf{V}(S_n)}{t^2} = \frac{\sum_{k=1}^n \sigma^2(X_k)}{t^2}.
\]

\end{thm}

\noindent Before we start the proof, note that the bound on the
right-hand side is the usual variance bound that
follows from Chebyshev's inequality; except that the
event on the left-hand side whose probability this
quantity bounds is a \emph{much bigger} event than the
usual deviation event $\{|S_n-\mathbf{E}(S_n)|>t\}$
for which we know the Chebyshev bound holds!

\begin{proof}
We may assume without loss of generality that
$\mathbf{E}(X_k)=0$ for all $k$. Denote

\[
A = \left\{ \max_{1\le k\le n} |S_k| > t\right\},
\]

\noindent and define events $A_1,A_2,\ldots,A_n$ by

\[
A_k = \left\{ |S_k|\ge t, \ \ \max_{1\le j<k} |S_k|<t \right\}.
\]

\noindent In words, $A_k$ is the event that the sequence of
cumulative sums $(S_j)_{j=1}^n$ exceeded $t$ in
absolute value for the first time at time $k$. Note
that these events are disjoint and their union is the
event $A$ whose probability we are trying to bound.
As a consequence, we can lower-bound the variance
$\mathbf{V}(S_n)=\mathbf{E}(S_n^2)$ of $S_n$, as
follows:

\begin{eqnarray*}
\mathbf{E}(S_n^2) &\ge& \mathbf{E}(S_n^2 1_A) = \sum_{k=1}^n \mathbf{E} \left(S_n^2 1_{A_k}\right)
= \sum_{k=1}^n \mathbf{E} \left[ \left(S_k+(S_n-S_k)\right)^2 1_{A_k} \right] \\
&=&
\sum_{k=1}^n \Big[\mathbf{E} (S_k^2 1_{A_k}) + \mathbf{E}\left[ (S_n-S_k)^2 1_{A_k} \right]
+ 2 \mathbf{E} \left[(S_k 1_{A_k}) (S_n-S-k) \right] \Big].
\end{eqnarray*}

\noindent In this last expression, the terms
$\mathbf{E} \left[(S_k 1_{A_k}) (S_n-S-k) \right]$
are equal to $0$, since $S_k 1_{A_k}$ is a random
variable that depends only on $X_1,\ldots,X_k$ and
hence independent of $S_n-S_k$, which depends only on
the values of $X_{k+1},\ldots,X_n$, which causes the
expectation of their product to be equal to the
product of the expectations, which is 0. Furthermore,
the middle terms
$\mathbf{E}\left[ (S_n-S_k)^2 1_{A_k} \right]$ are
all nonnegative, and each of the first terms
$\mathbf{E} (S_k^2 1_{A_k})$ satisfies

\[
\mathbf{E} (S_k^2 1_{A_k}) \ge \mathbf{E} (t^2 1_{A_k}) = t^2 \mathrm{P}(A_k),
\]

\noindent since on the event $A_k$ we know that $S_k^2$ is at
least $t^2$ (look again at the definition of $A_k$).
Combining these observations, we get that

\[
\mathbf{V}(S_n)  \ge t^2\sum_{k=1}^n \mathrm{P}(A_k) = t^2 \mathrm{P}(A),
\]

\noindent which is exactly the claim that was to be proved.
\end{proof}

\noindent As a corollary, we get a result on convergence of
random series.

\begin{thm}
\label{thm-one-series}
Let $X_1,X_2,\ldots$ be a sequence of independent
r.v.'s such that $\mathbf{E}(X_n)=0$ for all $n$, and
assume that
$\sum_{n=1}^\infty \mathbf{V}(X_n) < \infty$. Then
the random series $\sum_{n=1}^\infty X_n$ converges
almost surely.
\end{thm}

\begin{proof}
Denote as usual $S_n = \sum_{k=1}^n X_k$. We have the
following equality of events:

\begin{eqnarray*}
\left\{ \sum_{n=1}^\infty X_n\textrm{ converges}
\right\}
&=&
\left\{ \left(\sum_{n=1}^N X_n\right)_{N\ge 1} \textrm{ is a Cauchy sequence}
\right\} \\
&=&
\bigcap_{\epsilon>0}
\bigcup_{N\ge 1}
\bigcap_{n\ge N}
\Big\{
|S_n-S_N|<\epsilon
\Big\}.
\end{eqnarray*}

\noindent Or, put differently, we can look at the complement of
this event and represent it as

\begin{eqnarray*}
\left\{ \sum_{n=1}^\infty X_n\textrm{ does not converge}
\right\}
&=&
\bigcup_{\epsilon>0}
\bigcap_{N\ge 1}
\bigcup_{n\ge N}
\Big\{
|S_n-S_N|\ge\epsilon
\Big\} \\
&=&
\bigcup_{\epsilon>0}
\bigcap_{N\ge 1}
\left\{ \sup_{n\ge N}
|S_n-S_N|\ge\epsilon
\right\}.
\end{eqnarray*}

\noindent This form is exactly suitable for an application of
the Kolmogorov's maximal inequality, except that here
we have an infinite sequence of partial sums instead
of a finite maximum. However, by the ``continuity from
below'' property of probability measures, we see that
it does not matter. More precisely, for any
$\epsilon>0$ and $N\ge 1$, we have

\begin{eqnarray*}
\mathrm{P}\left(
\sup_{n\ge N}
|S_n-S_N|\ge\epsilon
\right) &=& \lim_{M\to\infty} \mathrm{P}\left(
\sup_{N\le n\le M}
|S_n-S_N|\ge\epsilon
\right) \le \lim_{M\to\infty} \frac{\mathbf{V}(S_M-S_N)}{\epsilon^2} \\ &=& \frac{1}{\epsilon^2} \sum_{n=N}^\infty \mathbf{V}(X_n).
\end{eqnarray*}

\noindent Therefore also

\begin{eqnarray*}
\mathrm{P}\left( \bigcap_{N\ge 1} \left\{
\sup_{n\ge N}
|S_n-S_N|\ge\epsilon
\right\} \right) &\le& \inf_{N\ge 1}
\mathrm{P}\left(
\sup_{n\ge N}
|S_n-S_N|\ge\epsilon
\right) \\ &\le&
\inf_{N\ge 1} \frac{1}{\epsilon^2} \sum_{n=N}^\infty \mathbf{V}(X_n) = 0,
\end{eqnarray*}

\noindent because of our assumption that the sum of the
variances converges. Finally, this is true for all
$\epsilon>0$, so by the usual trick of replacing an
uncountably-infinite intersection by a countable one
(provided that the particular form of the event in
question warrants this!), we get the claim that

\[
\mathrm{P}\left(\sum_{n=1}^\infty X_n\textrm{ does not converge}\right) = 0.
\]

\end{proof}

\noindent One could ask whether the sufficient condition given
by the theorem above is also necessary (it is). More
generally, what happens for random variables with
non-zero expectations? What happens for r.v.'s with
infinite expectations, or infinite variances?
Kolmogorov formulated a general theorem that gives a
necessary and sufficient condition for a series of
independent random variables to converge almost
surely.

\begin{thm}[The Kolmogorov three-series theorem]
\label{thm-three-series}
If $X_1,X_2,\ldots$ is a sequence of independent
random variables, then the random series
$\sum_{n=1}^\infty X_n$ converges almost surely if
and only if the following three conditions hold:

\begin{enumerate}[label=\arabic*.]
\item $\sum_{n=1}^\infty \mathrm{P}(|X_n|>1) < \infty$.

\item The series
$\sum_{n=1}^\infty \mathbf{E}(X_n
\mathbf{1}_{\{|X_n|\le 1\}})$
converges.

\item $\sum_{n=1}^\infty \mathbf{V}(X_n \mathbf{1}_{\{|X_n|
\le 1\}}) < \infty$.
\end{enumerate}

\noindent If one of the conditions does not hold, then series
$\sum X_n$ diverges almost surely.
\end{thm}

\noindent We postpone the proof of \hyperref[thm-three-series]{Theorem~\ref*{thm-three-series}}
until later; we will prove it in Section
\ref{sec-k-3-series} as an application of a generalized
version of the central limit theorem. Note that since
the convergence of the series $\sum X_n$ is
equivalent to the convergence of $\sum a X_n$ for any
constant $a$, the value $1$ chosen for the truncation
of $X_n$ in the theorem is arbitrary and can be
replaced by any other constant.

\begin{example}
Let $(c_n)_{n=1}^\infty$ be a sequence of real
numbers. Consider the
$\textbf{series with random signs}$ associated with
the sequence $(c_n)$, which we denote
$\sum_{n=1}^\infty \pm c_n$, and which more precisely
represents the series

\[
\sum_{n=1}^\infty c_n X_n,
\]

\noindent where $X_1, X_2, \ldots$ is a sequence of i.i.d.
r.v.'s taking the values $-1,+1$ with respective
probabilities $1/2,1/2$. By \hyperref[thm-one-series]{Theorem~\ref*{thm-one-series}}
it follows that if $\sum_{n=1}^\infty c_n^2 < \infty$
then the series $\sum_{n=1}^\infty \pm c_n$ converges
almost surely. From \hyperref[thm-three-series]{Theorem~\ref*{thm-three-series}}, one
can check easily that this condition is also
necessary, in other words that the series with random
series converges a.s. if and only if the series of
squares $\sum c_n^2$ converges. Thus, for example,
the $\textbf{harmonic series with random signs}$$\sum \pm \frac{1}{n}$ converges a.s., but the
analogous series of square root reciprocals
$\sum_n \pm \frac{1}{\sqrt{n}}$ diverges a.s. Compare
this to the series with \emph{alternating} signs
$\sum_n \frac{(-1)^n}{n}$ and
$\sum_n \frac{(-1)^n}{\sqrt{n}}$, both of which are
known to converge! Try to develop your intuition by
thinking about the reasons why the behavior for series
with alternating signs and that for random signs is
not the same.
\end{example}
